% ============================================================
% packages.tex – Pakete, Bibliografie und Layout-Konfiguration
% ============================================================
% Bereits in main.tex geladen:
%   microtype, graphicx, fontspec, fancyhdr, setspace
% ============================================================


% ============================================================
% 1. Sprache, Zitate, Datum
% ============================================================
\usepackage[ngerman]{babel}        % Deutsche Sprache
\usepackage{csquotes}              % Korrekte Anführungszeichen, von biblatex benötigt
\usepackage[iso, german]{isodate}  % Deutsches Datumsformat (für \today)


% ============================================================
% 2. Seitenlayout: Spacing und Hilfsmittel
% ============================================================
\usepackage{parskip}
\setlength{\headheight}{15pt}
% Großer Stretch (plus80pt) ist Absicht: erlaubt LaTeX, Absatzabstände
% beim Float-Layout flexibel zu strecken, ohne den Text zu sperren.
\setlength{\parskip}{11pt plus80pt minus10pt}
% Abstand zwischen Haupttext und Fußnotenbereich
\setlength{\skip\footins}{1.5cm}

\usepackage{float}      % \begin{figure}[H] für exakte Positionierung
\usepackage{lscape}     % Querformat-Seiten (\begin{landscape})
\usepackage{blindtext}  % \blindtext{} für Layout-Tests
\usepackage{easy-todo}  % \todo{Text} für Inline-Notizen
\usepackage{refcount}   % \getpagerefnumber{} – wird in toc.tex verwendet


% ============================================================
% 3. Literaturverzeichnis (BibLaTeX + Biber)
% ============================================================
% style=ext-authoryear  Autor-Jahr-Stil mit Erweiterungen
% autocite=footnote     \autocite{} setzt automatisch eine Fußnote
% dashed=false          Wiederholte Autoren ausschreiben, nicht mit "—" abkürzen
% sorting=nyt           Sortierung: Name, Year, Title
% maxbibnames=99        bis zu 99 Autoren im Literaturverzeichnis
% maxcitenames=2        in der Fußnote max. 2 Autoren, danach "et al."
% mincitenames=1        in der Fußnote mind. 1 Autor
% uniquename=init       Vornamen-Initialen bei Namensgleichheit
% giveninits=true       Vornamen werden zu Initialen gekürzt
\usepackage[
  style=ext-authoryear,
  autocite=footnote,
  dashed=false,
  sorting=nyt,
  maxbibnames=99,
  maxcitenames=2,
  mincitenames=1,
  uniquename=init,
  giveninits=true,
]{biblatex}

% Klammer um die Jahreszahl in der Fußnote
\DeclareInnerCiteDelims{footcite}{\bibopenparen}{\bibcloseparen}

% Slash zwischen Autoren in der Fußnote
\renewcommand*{\multinamedelim}{\addslash}
\renewcommand*{\finalnamedelim}{\multinamedelim}

% Doppelpunkt nach Autorname im Literaturverzeichnis
\DeclareDelimFormat[bib]{nametitledelim}{\addcolon\space}

% URL-Präfix
\DeclareFieldFormat{url}{Verfügbar unter: \urlstyle{same}\url{#1}}

% Titel ohne Hervorhebung, Journaltitel kursiv
\DeclareFieldFormat{title}{{#1}}
\DeclareFieldFormat{journaltitle}{\mkbibemph{#1}}

% Sortierung nach "Nachname, Vorname"
\DeclareNameAlias{sortname}{family-given}

% "Jg." statt Standard-Präfix für Journal-Volumen / Heftnummer
\DeclareFieldFormat[article]{volume}{\bibstring{jourvol}\addnbspace #1}
\DeclareFieldFormat[article]{number}{\bibstring{number}\addnbspace #1}

% Deutsche Strings anpassen
\DefineBibliographyStrings{ngerman}{
  jourvol   = {Jg\adddot},
  andothers = {{et al\adddot}},
}

% Abstand zwischen Bibliographie-Einträgen
\setlength{\bibitemsep}{9pt}

% Optional: Alle Quellen anzeigen, auch nicht zitierte
% \nocite{*}

% Einbinden der Literaturdateien
% ============================================================
% Literatur-Dateien einbinden
% ============================================================
% Empfehlung: Eine .bib-Datei pro Kapitel anlegen und hier einbinden.
% Export aus Zotero (Better BibTeX), JabRef oder Citavi.
%
% Verwendung im Text:
%   \autocite[Vgl.][S. 42]{schluessel}          → Fußnotenzitat mit Vgl.
%   \autocite[][S. 42]{schluessel}               → Fußnotenzitat ohne Vgl.
%   \autocites[Vgl.][S. 1]{q1}[Vgl.][S. 5]{q2}  → mehrere Quellen
% ============================================================

% Hauptdatei mit allen Quellen:
\addbibresource{bibliographies/resources/literatur.bib}

% Optional: kapitelweise aufteilen
% \addbibresource{bibliographies/resources/01-einleitung/01-einleitung.bib}
% \addbibresource{bibliographies/resources/02-theorie/02-01-theorie-unterkapitel.bib}
% \addbibresource{bibliographies/resources/03-methode/03-01-methode-unterkapitel.bib}
% \addbibresource{bibliographies/resources/04-ergebnisse/04-01-ergebnisse-unterkapitel.bib}
% \addbibresource{bibliographies/resources/05-diskussion/05-01-diskussion-unterkapitel}


% ============================================================
% 4. Tabellen
% ============================================================
\usepackage{booktabs}
\usepackage{longtable}  % Tabellen über mehrere Seiten
\usepackage{multirow}   % Verbundene Zeilen in Tabellen
\usepackage{array}
\renewcommand{\arraystretch}{1.5}  % Erhöht Zeilenabstand in Tabellen
\newcolumntype{P}[1]{>{\RaggedRight\arraybackslash}p{#1}}  % linksbündige p-Spalte


% ============================================================
% 5. URLs und Hyperlinks
% ============================================================
\usepackage{xurl}              % URL-Umbruch an beliebigen Stellen
\urlstyle{same}                % URLs im Stil der Grundschrift
\usepackage[hidelinks]{hyperref}  % Klickbare, aber unsichtbare Links


% ============================================================
% 6. Seitenränder
% ============================================================
% SRH-Vorgabe: 4 cm links (Korrekturrand), 2 cm rechts, 2,5/2 cm oben/unten
\usepackage[right=2cm, left=4cm, top=2.5cm, bottom=2cm]{geometry}


% ============================================================
% 7. Aufzählungen, Mathematik, Boxen
% ============================================================
\usepackage{enumitem}
\usepackage{ragged2e}
\usepackage{amsmath}
\usepackage{framed}    % \begin{framed} – Rahmen um Textblöcke

% Engere itemize-Variante mit minimalem Abstand
\newenvironment{tightitemize}
  {\begin{itemize}\setlength{\itemsep}{0pt}\setlength{\parskip}{0pt}\setlength{\parsep}{0pt}}
  {\end{itemize}}


% ============================================================
% 8. Abkürzungsverzeichnis
% ============================================================
% nohyperlinks   keine Sprunglinks im Text
% printonlyused  nur tatsächlich verwendete Abkürzungen drucken
\usepackage[nohyperlinks, printonlyused]{acronym}


% ============================================================
% 9. Bildunterschriften (captions)
% ============================================================
% raggedright + singlelinecheck=false: Captions immer linksbündig,
% auch wenn sie nur eine Zeile lang sind.
% font=small, format=hang: kleine Schrift, Folgezeilen unter den Text eingerückt.
\usepackage[
  justification=raggedright,
  singlelinecheck=false,
  font=small,
  format=hang,
]{caption}

% Trennzeichen "Tabelle 1:" mit Doppelpunkt + Tabulator-Abstand
\DeclareCaptionLabelSeparator{colonquad}{: \quad}
\captionsetup{labelsep=colonquad}

% Caption für longtable ebenfalls linksbündig
\captionsetup[longtable]{justification=raggedright, singlelinecheck=false}

% longtable-Kasten und Caption am linken Textblock ausrichten
\setlength{\LTleft}{0pt}
\setlength{\LTright}{0pt}
\setlength{\LTcapwidth}{\textwidth}


% ============================================================
% 10. Anhang
% ============================================================
\usepackage{appendix}
\usepackage{etoolbox}  % wird für \apptocmd im Anhang verwendet


% ============================================================
% 11. Inhaltsverzeichnis (TOC, LOF, LOT)
% ============================================================
\usepackage{etoc}
\renewcommand{\etocaftertitlehook}{\pagestyle{fancy}}
\renewcommand{\etocaftertochook}{\thispagestyle{fancy}}

% Seitenstil für Verzeichnis-Seiten: Seitenzahl mittig, keine Trennlinie
\fancypagestyle{plain}{
  \fancyhf{}
  \fancyhead[C]{\thepage}
  \renewcommand{\headrulewidth}{0pt}
  \renewcommand{\footrulewidth}{0pt}
}

\usepackage{tocloft}

% Präfix "Tabelle X:" im Tabellenverzeichnis
\renewcommand{\cfttabpresnum}{Tabelle }
\renewcommand{\cfttabaftersnum}{:}
\renewcommand{\cfttabnumwidth}{6.5em}
\setlength{\cfttabindent}{0pt}

% Präfix "Abbildung X:" im Abbildungsverzeichnis
\renewcommand{\cftfigpresnum}{Abbildung }
\renewcommand{\cftfigaftersnum}{:}
\renewcommand{\cftfignumwidth}{8em}
\setlength{\cftfigindent}{0pt}


% ============================================================
% 12. Anlagenverzeichnis (eigenes \listof)
% ============================================================
% Definiert ein neues Verzeichnis "Anlage" mit eigener Nummerierung.
% Verwendung im Anhang: \Anlage{Titel} → fügt einen Eintrag hinzu.
% Wichtig: \newlistof benötigt einen LEEREN Titel-Parameter, weil die
% Überschrift separat per \section* aus anlagenverzeichnis.tex kommt.
\newlistof{Anlage}{ana}{}

% Nummerierung arabisch (1, 2, 3, …)
\renewcommand*{\theAnlage}{\arabic{Anlage}}

% Darstellung im Verzeichnis: Prefix "Anhang", Doppelpunkt, Spacing wie Tab/Abb
\renewcommand{\cftAnlagepresnum}{Anhang }
\renewcommand{\cftAnlageaftersnum}{:}
\renewcommand{\cftAnlagenumwidth}{6.5em}
\setlength{\cftAnlageindent}{0pt}
\renewcommand{\cftAnlagedotsep}{\cftdotsep}
\renewcommand{\cftAnlagefont}{\normalfont}
\renewcommand{\cftAnlagepagefont}{\normalfont}

% Befehl zum Hinzufügen eines Eintrags
\newcommand{\Anlage}[1]{%
  \refstepcounter{Anlage}%
  \addcontentsline{ana}{Anlage}{\protect\numberline{\theAnlage}#1}%
}
